\chapter{Fusion of several consistency proofs}

\section{The infinitary system}
We adapt the system presented as $\textsf{IHA}$ in \cite{TAT11} to align with the finitary calculus:

\begin{cthdefinition}{(The system $\HA_\omega$)}
We define a sequent $\seqd{\Gamma}{C}{\alpha}{r}$ where $\Gamma$ is a multiset, $C$ is a formula, $r \in \mathbb{N}$ and $\alpha$ an ordinal. \emph{$\Gamma \vdash^\alpha_r C$ has a derivation with maximum height $\alpha$ and with with maximum cut $r$} if it is obtained from the inference rules on the \emph{initial sequents}. In the following, $\alpha_i < \alpha$ and $C$ may be empty.
\begin{itemize}
\item \textbf{\emph{Initial sequents:}}
\begin{align*}
&\overline{\seqd{\Gamma}{P}{\alpha}{r}}\;(\rm{Ax\;R}) \text{ for true atoms $P$ and } \\
&\overline{\seqd{\Gamma, P}{C}{\alpha}{r}}\;(\rm{Ax\;L}) \text{ for false atoms $P$}
\end{align*}

\item \textbf{\emph{Logical rules:}}
{\small
% AND Right/Left
\begin{center}
\AxiomC{$\seqd{\Gamma}{A}{\alpha_1}{r}$}
\AxiomC{$\seqd{\Gamma}{B}{\alpha_2}{r}$}
\RightLabel{$(\land\rm{R})$}
\BinaryInfC{$\seqd{\Gamma}{A\land B}{\alpha}{r}$}
\DisplayProof
\quad
\AxiomC{$\seqd{\Gamma,A,B}{C}{\alpha_1}{r}$}
\RightLabel{$(\land\rm{L})$}
\UnaryInfC{$\seqd{A\land B,\Gamma}{C}{\alpha}{r}$}
\DisplayProof
\end{center}

% OR Right/Left
\begin{center}
\AxiomC{$\seqd{\Gamma}{A_i}{\alpha_1}{r}$}
\RightLabel{$(\lor\rm{R})$}
\UnaryInfC{$\seqd{\Gamma}{A_0 \lor A_1}{\alpha}{r}$}
\DisplayProof
\quad
\AxiomC{$\seqd{\Gamma,A}{C}{\alpha_1}{r}$}
\AxiomC{$\seqd{\Gamma,B}{C}{\alpha_2}{r}$}
\RightLabel{$(\lor\rm{L})$}
\BinaryInfC{$\seqd{A\lor B,\Gamma}{C}{\alpha}{r}$}
\DisplayProof
\end{center}

% IMP Right/Left
\begin{center}
\AxiomC{$\seqd{\Gamma, A}{B}{\alpha_1}{r}$}
\RightLabel{$(\rightarrow\rm{R})$}
\UnaryInfC{$\seqd{\Gamma}{A\rightarrow B}{\alpha}{r}$}
\DisplayProof
\quad
\AxiomC{$\seqd{\Gamma,A\rightarrow B}{A}{\alpha_1}{r}$}
\AxiomC{$\seqd{B,\Gamma}{C}{\alpha_2}{r}$}
\RightLabel{$(\rightarrow\rm{L})$}
\BinaryInfC{$\seqd{\Gamma,A\rightarrow B}{C}{\alpha}{r}$}
\DisplayProof
\end{center}

% FORALL Right/Left
\begin{center}
\AxiomC{$\seqd{\Gamma}{A(t)}{\alpha_i}{r} \quad \text{(for all $t$)}$}
\RightLabel{$(\forall\rm{R}_\infty)$}
\UnaryInfC{$\seqd{\Gamma}{\forall x\,A(x)}{\alpha}{r}$}
\DisplayProof
\quad
\AxiomC{$\seqd{\Gamma,A(t),\forall x\,A(x)}{C}{\alpha_i}{r}$}
\RightLabel{$(\forall\rm{L})$}
\UnaryInfC{$\seqd{\Gamma,\forall x\,A(x)}{C}{\alpha}{r}$}
\DisplayProof
\end{center}

% EXISTS Right/Left
\begin{center}
\AxiomC{$\seqd{\Gamma}{A(t)}{\alpha_1}{r}$}
\RightLabel{$(\exists\rm{R})$}
\UnaryInfC{$\seqd{\Gamma}{\exists x\,A(x)}{\alpha}{r}$}
\DisplayProof
\quad
\AxiomC{$\Gamma,\seqd{A(t)}{C}{\alpha_i}{r} \quad \text{(for all $t$)}$}
\RightLabel{$(\exists\rm{L}_\infty)$}
\UnaryInfC{$\seqd{\Gamma,\exists x\,A(x)}{C}{\alpha}{r}$}
\DisplayProof
\end{center}}
\item \textbf{\emph{Structural rule:}}
\begin{center}
\AxiomC{$\seqd{\Gamma}{A}{\alpha_1}{r}$}
\AxiomC{$\seqd{\Gamma,A}{C}{\alpha_2}{r}$}
\RightLabel{$(\rm{Cut}) \quad \text{ where $r < \gr(A)$}$}
\BinaryInfC{$\seqd{\Gamma}{C}{\alpha}{r}$}
\DisplayProof
\end{center}
\end{itemize}
A derivation $\seqd{\Gamma}{C}{\alpha}{r}$ is cut-free if $r = 0$.


\end{cthdefinition}

The main properties we need to prove for $\HA_\omega$:

\begin{cthlemma}{(Weakening)}
If $\seqd{\Gamma}{C}{\alpha}{r}$, then $\seqd{\Gamma,\Pi}{C}{\alpha}{r}$.
\end{cthlemma}

\begin{cthlemma}{(Inversion)}
    The following holds for some $\alpha$ and $r$:
    \begin{enumerate}
    \item $\seqd{A\land B,\Gamma}{C}{\alpha}{r}$ implies
      $\seqd{A,B,\Gamma}{C}{\alpha}{r}$;
    \item $\seqd{A\lor B,\Gamma}{C}{\alpha}{r}$ implies
      $\seqd{A,\Gamma}{C}{\alpha}{r}$ and $\seqd{B,\Gamma}{C}{\alpha}{r}$;
    \item $\seqd{A\rightarrow B,\Gamma}{C}{\alpha}{r}$ implies
      $\seqd{B,\Gamma}{C}{\alpha}{r}$;
    \item $\seqd{\exists x A(x),\Gamma}{C}{\alpha}{r}$ implies
      $\seqd{\Gamma,A(t)}{C}{\alpha}{r}$ for every $t$;
    \item $\seqd{\Gamma}{A\land B}{\alpha}{r}$ implies
      $\seqd{\Gamma}{A}{\alpha}{r}$ and $\seqd{\Gamma}{B}{\alpha}{r}$;
    \item $\seqd{\Gamma}{A\rightarrow B}{\alpha}{r}$ implies
      $\seqd{A,\Gamma}{B}{\alpha}{r}$;
    \item $\seqd{\Gamma}{\forall x A(x)}{\alpha}{r}$ implies
      $\seqd{\Gamma}{A(n)}{\alpha}{r}$ for every $t$.
    \end{enumerate}
    $(\lor R)$, $(\exists R)$, $(\forall L)$, and $(\to L)$ are not entirely invertible.
\end{cthlemma}

\begin{cththeorem}{(Monotonicity)}
    If $\seqd{\Gamma}{C}{\alpha}{r}$, $\alpha < \beta$ and $r < s$, then $\seqd{\Gamma}{C}{\beta}{s}$
\end{cththeorem}

\section{Embedding}

To prove that $\HA$ embeds into $\HA_\omega$, we will want to prove induction admissibility:

\begin{cthlemma}{(Induction admissibility)}
    Given $\seqd{\Gamma}{A(0)}{\alpha}{r}$ and $\seqd{\Gamma,A(x)}{A(S(x))}{\beta}{r}$, then:
    \[ \seqd{\Gamma}{\forall x\; A(x)}{\gamma}{r}\]
\end{cthlemma}

However, this is not possible if our $\omega$-rule requires every term. As such we need to normalise our terms:

\begin{cthdefinition}{(Normalised terms)}
    For a closed formula $A$, we will call the normalised formula $A^*$, the result of replacing every closed term $t$ with $\val(t)$.
\end{cthdefinition}
Evaluating all the closed terms will not change the truth value of a sequent. For free variables, we will effectively move the free variables out of the object language and consider their universal closure. This is equivalent to replacing the free variables with every possible numeral:
\begin{align*}
    \seq{\Gamma}{a(a+b)=a^2+ab} \rightsquigarrow &\forall m,n \in \mathbb{N} \; \left( \seq{\Gamma^*}{\bar{m}(\bar{m}+\bar{n})=\bar{m}^2-\bar{n}} \right) \\
    \leftrightsquigarrow &\seq{(\Gamma[\sigma])^*}{(a(a+b)=a^2+ab)[\sigma]}
\end{align*}
\todo{Substitution notation}
Thus each sequent in this embedding will be mapped to multiple sequents. By making this accommodation, we find $\HA_\omega$ with only normalised formulae has admissibility of induction:

\begin{cthlemma}{(Induction admissibility)}
    Given $\seqd{\Gamma}{A(0)}{\alpha}{r}$ and $\forall n\in\mathbb{N}\;\left(\seqd{\Gamma,A(\nmrl{n})}{A(S(\nmrl{n}))}{\beta}{r}\right)$ where $\gr(A(x)) < r$ then:
    \[ \seqd{\Gamma}{\forall x\; A(x)}{\max(\alpha,\beta)+\omega}{r}\]
\end{cthlemma}
\begin{proof}
    Notice that we can combine instances of the premises with cut. As an example:
\begin{center}
\AxiomC{$\seqd{\Gamma}{A(0)}{\max(\alpha,\beta)+1}{r}$}
\AxiomC{$\seqd{\Gamma,A(0)}{A(\nmrl{1})}{\max(\alpha,\beta)+1}{r}$}
\RightLabel{$(\rm{Cut})$}
\BinaryInfC{$\seqd{\Gamma}{A(\nmrl{1})}{\max(\alpha,\beta)+1}{r}$}
\DisplayProof
\end{center}
    For each $n$, we will have to add another cut to obtain:
    \[ \seqd{\Gamma}{A(\nmrl{n})}{\max(\alpha,\beta)+n}{r}\]
    with $\gr(A(\nmrl{n})) < r$. Applying $\forall\rm{R}_\infty$, we obtain $\seqd{\Gamma}{\forall x\;A(x)}{\max(\alpha,\beta) + \omega}{r}$
\end{proof}

Induction will be the main contribution to an increase in height when translating a proof from $\HA$ to $\HA_\omega$. Thus each induction will add an extra $\omega$ to the height of the proof when converted to a $\HA_\omega$ proof:

\begin{cththeorem}{(Embedding)}
    If $\seqd{\Gamma}{C}{n}{r}$ is derivable in $\HA$ with $k$ inductions, then $\seqd{(\Gamma[\sigma])^*}{(C[\sigma])^*}{\omega \cdot k + n}{r}$ is derivable in $\HA_\omega$ where $\sigma$ is any substitution of free variables for numerals.
\end{cththeorem}
\begin{proof}
    By induction on the height of $\seqd{\Gamma}{C}{n}{r}$.
    \begin{itemize}
        \item \textbf{\emph{Initial sequents:}} If $P$ is an atom in $\seq{\Gamma,P}{P}$, then upon substitution $P[\sigma]$, we find it must become a closed atom. If it is true, then $(\rm{Ax\;R})$ obtains the desired sequent and $(\rm{Ax\;L})$ otherwise.
        \item \textbf{\emph{Arithmetical rules:}} \todo{Write up}
    \end{itemize}
\end{proof}
\section{Rank reduction}
Before proving rank reduction, we need the following lemma:
\begin{cthlemma}{(Single-cut Rank Reduction)}
    Let $\gr(A) = r$, $\seqd{\Gamma}{A}{\alpha}{r}$ and $\seqd{\Pi,A}{C}{\beta}{r}$, then:
    \[ \seqd{\Gamma,\Pi}{C}{\omega \cdot (\alpha \# \beta) + \omega}{r} \]
\end{cthlemma}
For any cut-formula $A$ in a proof $\seqd{\Gamma}{C}{\alpha}{r}$ we require $\gr(A) < r$. Thus this lemma requires replacing the application of cut on $A$ to one on its subformulae. As an example, if $A = B \to B'$, suppose $A$ is principal formula on both sides of the cut. Therefore we have:
\begin{center}
    \AxiomC{$\seqd{\Gamma}{B}{\beta'}{r}$}
    \AxiomC{$\seqd{\Gamma,B'}{C}{\beta'}{r}$}
    \RightLabel{$(\to\rm{R})$}
    \BinaryInfC{$\seqd{\Gamma,B \to B'}{C}{\beta'}{r}$}
    \DisplayProof
    \quad
    \AxiomC{$\seqd{\Gamma,B}{B'}{\alpha'}{r}$}
    \RightLabel{$(\to\rm{L})$}
    \UnaryInfC{$\seqd{\Gamma}{B \to B'}{\alpha'}{r}$}
    \DisplayProof
\end{center}
We can therefore, apply cut on the premises:
\begin{center}
    \AxiomC{$\seqd{\Gamma}{B}{\gamma}{r}$}
    \AxiomC{$\seqd{\Gamma,B}{B'}{\gamma}{r}$}
    \AxiomC{$\seqd{\Gamma,B'}{C}{\gamma}{r}$}
    \RightLabel{$(\rm{Cut})$}
    \BinaryInfC{$\seqd{\Gamma,B}{C}{\gamma}{r}$}
    \RightLabel{$(\rm{Cut})$}
    \BinaryInfC{$\seqd{\Gamma}{C}{\gamma}{r}$}
    \DisplayProof
\end{center}
Therefore, we see this subproof contributes at most $\max(\alpha',\beta')+2$. If we have $n$ many $A = B \to B'$ in the proof, we can bound the whole proof with $\omega \cdot (\alpha \# \beta) + \omega$ (since $\max(\alpha,\beta) < \alpha \# \beta$). 

\begin{cththeorem}{(Rank reduction)}
    If $\seqd{\Gamma}{C}{\alpha}{r+1}$, then $\seqd{\Gamma}{C}{\omega^{\left(\omega^\alpha\right)}}{r}$
\end{cththeorem}
\begin{proof}
    This is by induction on the left and right derivations. \todo{Write up the other cases (inversion lemmas helpful here)}. Let $A$ be the cut with $\gr(A) = r$, then we have premises:  $\seqd{\Gamma}{A}{\alpha_1}{r+1}$ and $\seqd{\Gamma,A}{C}{\alpha_2}{r+1}$ (where $\alpha_1,\alpha_2 < \alpha$). Applying the induction hypothesis on the two premises gives us: $\seqd{\Gamma}{A}{\omega^{\omega^\alpha_1}}{r}$ and $\seqd{\Gamma,A}{C}{\omega^{\omega^\alpha_2}}{r}$. We can apply our previous lemma on this to obtain: 
    \[ \seqd{\Gamma}{C}{\omega\cdot \left(\omega^{\omega^{\alpha_1}}\#\omega^{\omega^{\alpha_2}}\right) + \omega}{r} \]
    Since $\omega\cdot \left(\omega^{\omega^{\alpha_1}}\#\omega^{\omega^{\alpha_2}}\right) + \omega < \omega^{\left(\omega^\alpha\right)}$, our conclusion holds.
\end{proof}

\begin{cththeorem}{(Cut-elimination)}
    If $\seqd{\Gamma}{C}{\alpha}{r}$, then $\seqd{\Gamma}{C}{\omega_{2r}(\alpha)}{0}$
\end{cththeorem}
\begin{proof}
    By induction on $r$, an repeated application of Rank reduction. For the inductive step, we have $\seqd{\Gamma}{C}{\alpha}{r+1}$, then by Rank reduction $\seqd{\Gamma}{C}{\omega^{\left(\omega^\alpha\right)}}{r}$. This is equivalent to $\seqd{\Gamma}{C}{\omega_2(\alpha)}{r}$, which we apply our induction hypothesis will reduce to $\seqd{\Gamma}{C}{\omega_{2r}(\omega_2(\alpha))}{0}$. Finally this is equivalent to $\seqd{\Gamma}{C}{\omega_{2(r+1)}(\alpha)}{0}$
\end{proof}

\begin{cththeorem}{(No cut-free empty sequent)}
    $\seqd{}{}{\alpha}{r}$ does not hold.
\end{cththeorem}
\begin{proof}
    This holds by induction on number of inferences.
\end{proof}

\begin{cththeorem}{(Consistency of HA)}
    $\HA$ is consistent.
\end{cththeorem}

\todo{Currently $\HA_\omega$ proves $\HA$ can't, but if we can change that (i.e. weaken the $\omega$-rule), then I see nowhere the proof will go wrong. And if thats the case, it seems like I can reobtain disjunction/existential property for $\HA$ and $\HA_\omega$ }

%%% Local Variables:
%%% mode: LaTeX
%%% TeX-master: "../main"
%%% End:
